\documentclass[conference]{IEEEtran}

\usepackage{amsmath,amssymb,bm}
\usepackage{booktabs}
\usepackage{cite}
\usepackage{url}

\newcommand{\Ftwo}{\mathbb{F}_2}
\newcommand{\R}{\mathbb{R}}

% Title: CL-52c amended by CL-70, 2026-08-05.
\title{Differentiable Joint Estimation of the Partial-FFT Response,
Combiner, and Coded Bits}

\author{\IEEEauthorblockN{Gabriel Chua Yu Han}}

\begin{document}
\maketitle

\begin{abstract}
Underwater acoustic communication channels are time varying and frequency
selective.  Multipath and ambient noise make uncoded symbol decisions
unreliable, so transmissions carry forward error correction (FEC) that
couples every coded carrier in the codeword.  Partial fast Fourier transform
(Partial-FFT) demodulation is an established defense against residual Doppler,
and it keeps short-time observations of the symbol interval.  It has been used
ahead of the decoder: pilot-trained combiner weights form pre-decoder soft
symbols, and FEC decoding follows.  However, where
pilots are sparse and residual Doppler strong, separating demodulation from
decoding discards the information the code carries during the combiner fit.
We keep the same Partial-FFT demodulation, but formulate one objective over
the Partial-FFT response, combiner, and relaxed codeword bits; observations,
pilots, and parity constraints regularize one another.  This joint maximum a posteriori (MAP)
problem is intractable because bits are discrete while the response and
combiner are continuous.  We therefore relax coded bits to real variables and
replace parity checks with differentiable penalties.  We show that completeness
and exactness trade against each other; the fewer variables the gradient moves,
the more of the objective is solved in closed form.  The code acts during
symbol estimation rather than after it.
\end{abstract}

\begin{IEEEkeywords}
underwater acoustic communication, OFDM, Doppler, partial FFT, LDPC, belief
propagation, maximum a posteriori (MAP), joint estimation, iterative
receiver.
\end{IEEEkeywords}

\section{Introduction}

Underwater acoustic communication channels are time varying and frequency
selective.  Multipath and ambient noise make uncoded symbol decisions
unreliable, so transmissions carry forward error correction (FEC).  One
popular choice is a linear block code, whose codeword satisfies parity
constraints; among these, low-density parity-check (LDPC) codes are the most
widely used.  This paper is not limited to LDPC: any linear block code with
parity constraints serves, and we use an LDPC code as the demonstration.
The parity constraints couple every coded carrier in the codeword, and they
are the only constraints in the receiver that do so; pilots constrain their
own carriers and their neighbours.

Orthogonal frequency-division multiplexing (OFDM) permits frequency-domain
equalization when the channel is constant over one symbol interval.  Residual
Doppler removes that condition.  After initial resampling a residual Doppler
factor and a carrier frequency offset remain, the channel varies within the
interval, and the frequency-domain operator is no longer diagonal.  Its
off-diagonal entries are intercarrier interference (ICI) coefficients, and a
single full-interval fast Fourier transform (FFT) collapses the time-local
information that would separate them.

Partial-FFT demodulation divides the useful interval into segments and keeps
one observation from every segment~\cite{yerramalli}.  It estimates no ICI
coefficient.  It keeps time-local components of the same full-interval
statistic, and complex combiner weights recombine them into one pre-decoder
soft symbol per carrier.  The weights are fitted to pilots, and FEC decoding
follows.  While this receiver is strong and
costs only one small ridge fit per carrier band, it consults the code only
after the soft symbols are formed.

That separation discards information.  Each soft symbol is formed as though
the coded carriers were independent, although the parity constraints tie
them.  Where the pilots are dense enough to track the coupling, little is
lost.  However, where pilots are sparse and residual Doppler is strong, the
weights are fitted from a few nearby pilots.  The constraints that could have
shaped that fit are invoked only after it is complete.

Turbo equalization established the value of exchanging information between
adaptive equalization and channel decoding~\cite{laot,zheng}.  We apply that
principle to Partial-FFT reception and carry it one step further.  Rather
than pass information between two fitted stages, we place the Partial-FFT
response, the combiner weights, and the coded bits in a single
objective~\cite{juna}.  Observations, pilots, and parity constraints then
regularize one another while the symbols are still being estimated.  The
joint maximum a posteriori (MAP) detection problem over these three
quantities is intractable, since the bits are discrete while the response and
the combiner are continuous, and a discrete parity check admits no gradient.
We therefore relax each coded bit to a real variable and replace each parity
check by a differentiable penalty on the relaxed bits.  Every term of the
objective then has a gradient with respect to every quantity.

Two receivers follow from that objective, and they differ in whether the
combiner weights are estimated at all.  In the first, the weights are
computed from the estimated response of the wanted carrier, so the relaxed
bits are the only quantity searched.  In the second, the weights are
estimated together with the response and the bits, and a step is kept only
when the objective does not increase and the pilot error stays within a fixed
tolerance.  We compare the two on equal terms and do not argue that either is
the correct design.  The comparison is narrow by construction: the two
receivers share every other setting, so their decoded bit error rates measure
one difference alone.

We develop this in four steps.  First, we state the joint estimation and
decoding problem in its most general form and relax it.  Second, we write
the Partial-FFT observation model and the pilot-trained combiner that the
separate receiver fits.  Third, we specialize the objective and give the two
receivers: the response solved exactly whenever the bits are held fixed, and
the combiner either computed from that response or estimated beside it.
Fourth, we report what the comparison measures and the regime in which the
difference should appear.

\section{Joint Estimation and Decoding}

We wish to determine two unknowns from the received observations alone: the
coded bits, and the channel parameters through which those bits were
observed.  This section states that problem in its most general form, before
any receiver structure is chosen, and relaxes it so that descent can solve
it.

Let \(\mathbf b\in\Ftwo^{n_b}\) be the transmitted codeword of a linear
block code with parity-check matrix \(\mathbf P\in\Ftwo^{m\times n_b}\), and
let
\begin{equation}
  \mathcal C=\left\{\mathbf b\in\Ftwo^{n_b}:\mathbf P\mathbf b^{T}=\mathbf 0
  \ \text{over}\ \Ftwo\right\}
  \label{eq:codebook}
\end{equation}
be the codebook, where \(n_b\) is the coded-bit length and \(m\) is the
number of parity constraints.  Every transmission is a member of
\(\mathcal C\), so the code is a property of the transmitted signal itself,
not a processing stage of the receiver.

The modulation map turns the bits into transmitted symbols
\(\mathbf s(\mathbf b)\), and the channel turns the symbols into
observations:
\begin{equation}
  \mathbf Y=\mathcal A\!\left(\bm\psi,\mathbf s(\mathbf b)\right)+\mathbf n,
  \label{eq:general-forward}
\end{equation}
where \(\mathbf Y\) collects every observation the receiver keeps,
\(\mathcal A\) is the forward map, \(\bm\psi\) collects the unknown channel
parameters, and \(\mathbf n\) is noise.  Everything the receiver will ever
know about the bits and the channel is in \(\mathbf Y\), and both unknowns
enter it together.

The joint maximum a posteriori (MAP) formulation asks for both at once:
\begin{equation}
  (\widehat{\mathbf b},\widehat{\bm\psi})
  =
  \arg\max_{\mathbf b\in\mathcal C,\;\bm\psi}
  p(\mathbf Y\mid\mathbf b,\bm\psi)\,p(\bm\psi),
  \label{eq:general-map}
\end{equation}
where \(p(\mathbf Y\mid\mathbf b,\bm\psi)\) is the likelihood under
\eqref{eq:general-forward} and \(p(\bm\psi)\) carries what is known about
the channel before the observation, including the pilots.  Neither unknown
is conditioned on a fixed estimate of the other, so a constraint on either
one shapes the estimate of both.

A separate receiver factorizes \eqref{eq:general-map} instead.  It
estimates \(\bm\psi\) from pilots, forms soft symbols, and decodes
afterward.  While each factor is then a small problem, the bits are decided
against a channel estimate the code never shaped.

Two features make \eqref{eq:general-map} intractable.  The codebook doubles
with every message bit, so the search over \(\mathcal C\) cannot be carried
out directly.  The bits are discrete while \(\bm\psi\) is continuous, so no
gradient connects the two searches.

We therefore relax the bits.  Each coded bit takes a real variable
\(z_i\), and its relaxed bipolar value is
\begin{equation}
  \xi_i(\mathbf z)=\tanh z_i,
  \qquad z_i\in\R,
  \label{eq:relaxed-bit}
\end{equation}
where \(\xi_i=\pm1\) recovers the hard bit and intermediate values carry its
uncertainty.  The relaxed symbols \(\mathbf s(\mathbf z)\) follow through
the same modulation map.  The parity constraints become a penalty on the
relaxed bits:
\begin{equation}
  \mathcal R_{\mathrm{par}}(\mathbf z)
  =
  \sum_{j=1}^{m}
  \left(1-\prod_{i\in\mathcal N_j}\xi_i(\mathbf z)\right)^{2},
  \label{eq:parity-penalty}
\end{equation}
where \(\mathcal N_j\) is the set of bits that parity check \(j\)
constrains.  For hard bipolar values the product equals one exactly when the
check is satisfied, so the penalty vanishes exactly on codewords.

The relaxed joint problem is then one objective:
\begin{equation}
  \mathcal L(\bm\psi,\mathbf z)
  =
  -\log p\!\left(\mathbf Y\mid\mathbf s(\mathbf z),\bm\psi\right)
  -\log p(\bm\psi)
  +\lambda_{\mathrm p}\,\mathcal R_{\mathrm{par}}(\mathbf z),
  \label{eq:general-objective}
\end{equation}
where \(\lambda_{\mathrm p}\) weighs the parity penalty against the
likelihood.  Every term has a gradient with respect to every variable, so
demodulation and decoding can improve by the same descent.  The objective is
a surrogate, not the MAP problem itself: its minimizer depends on
\(\lambda_{\mathrm p}\), the relaxed bits need not be hard, and validity is
still checked on the hardened bits.

The next section fixes the forward map.  Partial-FFT reception determines
what \(\mathbf Y\) is, and the channel parameters \(\bm\psi\) split into the
response and the combiner weights.

\begin{thebibliography}{00}

\bibitem{yerramalli}
S. Yerramalli, M. Stojanovic, and U. Mitra,
``Partial FFT Demodulation: A Detection Method for Highly Doppler Distorted
OFDM Systems,''
\emph{IEEE Transactions on Signal Processing}, 2012.

\bibitem{laot}
C. Laot, A. Glavieux, and J. Labat,
``Turbo Equalization: Adaptive Equalization and Channel Decoding Jointly
Optimized,''
\emph{IEEE Transactions on Communications}, 2001.

\bibitem{zheng}
Y. R. Zheng, J. Wu, and C. Xiao,
``Turbo Equalization for Single-Carrier Underwater Acoustic Communications,''
\emph{IEEE Communications Magazine}.

\bibitem{juna}
G. C. Y. Han,
\emph{JUNA: A Joint Receiver for Doppler-Resilient OFDM: Partial-FFT
Coefficient Estimation and LDPC-Coupled Detection}.

\end{thebibliography}

\end{document}
